% Go Basics cheat sheet -- 2 pages, A4, 2 columns
% Compile with: pdflatex go_cheatsheet.tex
\documentclass[9pt]{extarticle}

\usepackage[a4paper, margin=1cm, top=0.9cm, bottom=1.35cm, footskip=0.6cm]{geometry}
\usepackage[T1]{fontenc}
\usepackage{textcomp}
\usepackage[utf8]{inputenc}
\usepackage{lmodern}
\usepackage[scaled=0.92]{helvet}
\renewcommand{\familydefault}{\sfdefault}
\usepackage[table]{xcolor}
\usepackage{multicol}
\usepackage{listings}
\usepackage{enumitem}
\usepackage{titlesec}
\usepackage{tabularx}
\usepackage{booktabs}
\usepackage{needspace}
\usepackage{etoolbox}
\pretocmd{\section}{\needspace{6\baselineskip}}{}{}

% ---------- colours (same as the slides) ----------
\definecolor{accent}{HTML}{204A96}
\definecolor{goblue}{HTML}{00ADD8}
\definecolor{ink}{HTML}{14213D}
\definecolor{body}{HTML}{3D4A5C}
\definecolor{codebg}{HTML}{EEF1F5}
\definecolor{comment}{HTML}{6B7F94}
\definecolor{rule}{HTML}{DDE1DA}
\definecolor{tablehead}{HTML}{E4E9E6}

\newcommand{\Go}{\textbf{\textcolor{goblue}{Go}}}
\newcommand{\kw}[1]{\textcolor{accent}{\texttt{#1}}}

\color{body}
\setlength{\parindent}{0pt}
\setlength{\parskip}{1pt}
\setlength{\columnsep}{0.6cm}
\setlength{\columnseprule}{0.3pt}
\renewcommand{\columnseprulecolor}{\color{rule}}

% ---------- headings ----------
\titleformat{\section}{\color{accent}\normalsize\bfseries}{}{0pt}{\MakeUppercase}[{\color{rule}\titlerule}]
\titlespacing*{\section}{0pt}{4pt}{1.5pt}
\titleformat{\subsection}{\color{ink}\small\bfseries}{}{0pt}{}
\titlespacing*{\subsection}{0pt}{3pt}{1pt}

% ---------- lists ----------
\setlist[itemize]{leftmargin=9pt, itemsep=0pt, parsep=0pt, topsep=1pt, label={\textcolor{accent}{\textbullet}}}

% ---------- code ----------
\lstdefinelanguage{Go}{
  morekeywords={break,case,chan,const,continue,default,defer,else,fallthrough,
    for,func,go,goto,if,import,interface,map,package,range,return,select,
    struct,switch,type,var,nil,true,false,iota},
  sensitive=true,
  morecomment=[l]{//},
  morecomment=[s]{/*}{*/},
  morestring=[b]",
  morestring=[b]`,
}
\lstset{
  language=Go,
  basicstyle=\ttfamily\footnotesize\color{ink},
  keywordstyle=\color{accent}\bfseries,
  commentstyle=\color{comment},
  stringstyle=\color{ink},
  backgroundcolor=\color{codebg},
  frame=none,
  xleftmargin=3pt, xrightmargin=3pt,
  framexleftmargin=3pt, framexrightmargin=3pt,
  aboveskip=2pt, belowskip=2pt,
  columns=fullflexible, keepspaces=true,
  showstringspaces=false,
  upquote=true,
  extendedchars=true,
  literate={é}{{\'e}}1,
}

\usepackage{fancyhdr}
\pagestyle{fancy}
\fancyhf{}
\renewcommand{\headrulewidth}{0pt}
\fancyfoot[L]{\footnotesize\color{body}Shane Petersen \textperiodcentered{} github.com/Mayoyama \textperiodcentered{} CC BY 4.0}
\fancyfoot[R]{\footnotesize\color{body}\thepage}

\begin{document}

% ---------- title ----------
{\color{ink}\LARGE\bfseries \Go{} Basics Cheat Sheet}\hfill
{\small For Python and C/C++ coders starting their first \Go{} project}\par
\vspace{1pt}{\color{rule}\rule{\linewidth}{0.6pt}}
\vspace{-4pt}

\begin{multicols*}{2}
\raggedright

% ======================================================
\section{What is Go}
\begin{itemize}
  \item Compiled like C; syntax reads somewhere between plain C and Python
  \item Garbage collected: no manual \kw{free} / \kw{delete}
  \item Built-in concurrency: goroutines and channels (no manual threads)
  \item No classes or inheritance; interfaces give duck-typing-like behaviour
  \item Generics since \Go{} 1.18
\end{itemize}
\subsection{Strengths}
\begin{itemize}
  \item Fast compiles; a single binary, no runtime to install
  \item Small language (25 keywords), one style enforced by \kw{gofmt}
  \item Strong standard library: networking, HTTP, JSON, files
\end{itemize}
\subsection{Weaknesses}
\begin{itemize}
  \item Verbose error handling (\kw{if err != nil} everywhere)
  \item Less memory control than C/C++/Rust, and usually slower
  \item Fewer features by design: no exceptions, no inheritance
  \item Nil pointer crashes are still possible at runtime
\end{itemize}

% ======================================================
\section{Running code}
\begin{lstlisting}
package main

import "fmt"

func main() {
    fmt.Println("hello")
}
\end{lstlisting}
\begin{itemize}
  \item \kw{go run main.go}: compile and run, binary is thrown away
  \item \kw{go build -o hello}: compile and keep the binary
  \item \kw{go vet ./...}: check for suspicious mistakes (\kw{./...} = this folder and every folder below)
  \item \kw{gofmt file.go} prints only; \kw{go fmt ./...} or \kw{gofmt -w .} rewrites files
  \item \kw{golangci-lint run ./...}: stricter, bundles many linters (separate install)
\end{itemize}

% ======================================================
\section{Packages \& imports}
\begin{lstlisting}
import (
    js "encoding/json"   // alias, like `as js`
    "unicode/utf8"       // full path to ONE package
)
n := utf8.RuneCountInString("Intéressant!")
data, _ := js.Marshal(n)   // use the last part / alias
\end{lstlisting}
\begin{itemize}
  \item A package = a folder of \kw{.go} files, each starting with \kw{package <name>}
  \item A program needs one \kw{package main} with one \kw{func main()}
  \item Unused imports \textbf{and} unused variables are compile errors
\end{itemize}

% ======================================================
\section{Naming \& visibility}
\begin{lstlisting}
func ParseInput() {}   // public: other packages can call it
func parseInput() {}   // private to this package
var MaxUsers = 10      // public
var userCount = 0      // private
\end{lstlisting}
\begin{itemize}
  \item camelCase / PascalCase; snake\_case only for file names
  \item \textbf{Uppercase} first letter = public (exported); \textbf{lowercase} = private to the package
  \item Applies to functions, types, struct fields and variables
  \item Enforced by the compiler (Python's \kw{\_private} is only a convention)
  \item The boundary is the package, not the struct: no class-level privacy
\end{itemize}

% ======================================================
\section{Variables \& zero values}
\begin{lstlisting}
var a int            // 0
var b int = 5
g := 5               // type inferred, functions only
const maxSize = 100  // compile-time, basic types only

var s string         // ""
var ok bool          // false
var p *int           // nil
var nums []int       // nil, len 0: append still works
var ages map[string]int  // nil: reading ok, adding panics
\end{lstlisting}
\begin{itemize}
  \item Every variable starts at its zero value, never garbage like in C
\end{itemize}

% ======================================================
\section{Numbers \& bytes}
\begin{itemize}
  \item \kw{int} (32/64 bit, default), \kw{int8}..\kw{int64}, \kw{uint8}..\kw{uint64}, \kw{byte}=\kw{uint8}
  \item \kw{float64} (default), \kw{float32}, \kw{bool}
  \item \kw{7 / 2} is \kw{3} (truncates like C); no \kw{//} operator (that's a comment)
  \item Overflow wraps silently: \kw{int8} 127 + 1 = -128 (\kw{math.MaxInt8}, ...)
  \item \kw{[]byte}: raw bytes for files and network; \kw{[]byte(s)} / \kw{string(b)}
\end{itemize}

% ======================================================
\section{Enums}
\begin{lstlisting}
type Color int
const (
    Red Color = iota   // 0, then counts up
    Green              // 1
    Blue               // 2
)
\end{lstlisting}
\begin{itemize}
  \item No \kw{enum} keyword: a custom int type plus a \kw{const} block
  \item No value check: \kw{Color(7)} compiles, while Python's \kw{Color(7)} raises \kw{ValueError}
  \item C++ equivalent: \kw{enum class Color \{ Red, Green, Blue \};}
\end{itemize}

% ======================================================
\section{Strings \& runes}
\begin{lstlisting}
s := "héllo"
len(s)                       // 6 bytes, not 5
utf8.RuneCountInString(s)    // 5 runes
s = "H" + s[1:]              // immutable: new string
raw := `no \n escapes`       // like Python r"..."
msg := fmt.Sprintf("hi %s, you are %d", "Ana", 30)
m := fmt.Sprintln("hi Ana", 30)  // "hi Ana 30\n"
b := []byte(s)
b[0] = 'J'
s = string(b)                // "Jéllo"
r := 'c' + 2                 // rune 101; string(r) is "e"
r = 'c' + rune(n)            // int var needs conversion
for i, ch := range s {}      // byte index, rune
\end{lstlisting}
\begin{itemize}
  \item \kw{"text"} is a string, \kw{\textquotesingle a\textquotesingle} is a rune (= \kw{int32}); no char type
  \item Strings are UTF-8 bytes: \kw{len(s)} counts bytes; \kw{"c" + 2} won't compile
  \item Backtick raw strings have no escapes and can span several lines
  \item \kw{Sprint}/\kw{Sprintf}/\kw{Sprintln} return a string instead of printing (like Python f-strings)
\end{itemize}

% ======================================================
\section{Slices}
\renewcommand{\arraystretch}{1.1}
{\footnotesize\ttfamily
\begin{tabularx}{\linewidth}{@{}>{\raggedright\arraybackslash}p{0.33\linewidth}X@{}}
\rowcolor{tablehead}\textsf{\textbf{Python}} & \textsf{\textbf{\Go}}\\
nums = [1, 2, 3]  & nums := []int\{1, 2, 3\}\\
nums.append(4)    & nums = append(nums, 4)\\
nums[1:3]         & nums[1:3]\\
nums[-1]          & nums[len(nums)-1]\\
3 in nums         & slices.Contains(nums, 3)\\
nums.sort()       & slices.Sort(nums)\\
\end{tabularx}}
\begin{itemize}
  \item \kw{append} is built in (no import); \kw{in} isn't a keyword
  \item \kw{slices} (1.21+): \kw{Max}, \kw{Min}, \kw{Reverse}, \kw{Equal}, \kw{Insert}, \kw{Delete}; \kw{sort.Ints}, \kw{sort.Strings}
  \item \kw{strings}: \kw{ToUpper}, \kw{ToLower}, \kw{TrimSpace}, \kw{ReplaceAll}, \kw{HasPrefix}, \kw{HasSuffix}, \kw{Split}, \kw{Join}
  \item \kw{[3]int} is a fixed-size array like C; slices are what you'll mostly use
\end{itemize}

% ======================================================
\section{Maps}
{\footnotesize\ttfamily
\begin{tabularx}{\linewidth}{@{}>{\raggedright\arraybackslash}p{0.33\linewidth}X@{}}
\rowcolor{tablehead}\textsf{\textbf{Python}} & \textsf{\textbf{\Go}}\\
ages = \{\}          & ages := map[string]int\{\}\\
ages["ana"] = 25    & ages["ana"] = 25\\
del ages["bob"]     & delete(ages, "bob")\\
"bob" in ages       & \_, ok := ages["bob"]\\
\end{tabularx}}
\begin{itemize}
  \item A missing key gives the zero value, so check \kw{ok}
  \item Iteration order is random (Python dicts keep insertion order)
\end{itemize}

\subsection{C++98 equivalents}
{\footnotesize\ttfamily
\begin{tabularx}{\linewidth}{@{}>{\raggedright\arraybackslash}p{0.36\linewidth}>{\raggedright\arraybackslash}X@{}}
\rowcolor{tablehead}\textsf{\textbf{\Go}} & \textsf{\textbf{C++98}}\\
append(nums, 4)          & nums.push\_back(4);\\
nums[1:3]                & std::vector<int>(nums.begin() + 1, nums.begin() + 3)\\
nums[len(nums)-1]        & nums.back()\\
slices.Contains(nums, 3) & std::find(nums.begin(), nums.end(), 3) != nums.end()\\
slices.Sort(nums)        & std::sort(nums.begin(), nums.end());\\
delete(ages, "bob")      & ages.erase("bob");\\
\_, ok := ages["bob"]    & ages.count("bob")\\
\end{tabularx}}
\begin{itemize}
  \item \kw{find}/\kw{sort} need \kw{<algorithm>}; \kw{std::map} keeps keys sorted
\end{itemize}

% ======================================================
\section{Errors are values}
\begin{lstlisting}
_, err := fmt.Println("hi")               // even Println returns one
err = errors.New("file is empty")         // fixed message
err = fmt.Errorf("cannot open %s", name)  // with values
\end{lstlisting}
\begin{itemize}
  \item \kw{error} is a type, returned like any other value; \kw{nil} means no error
  \item \kw{errors.New} needs \kw{import "errors"}; \kw{fmt.Errorf} formats like \kw{Sprintf}
  \item Messages: lowercase, no ending punctuation
\end{itemize}
\subsection{Two ways to check}
\begin{lstlisting}
n, err := strconv.Atoi("42")   // 1. check on the next line
if err != nil {
    return err
}
fmt.Println(n)                 // n exists here

if x, ok := ages["bob"]; !ok { // 2. check inside the if
    fmt.Println("not found")
} else {
    fmt.Println(x)
}                              // x does not exist here
\end{lstlisting}
\begin{itemize}
  \item \kw{ok} is a bool (map lookups, type assertions, channel receives); \kw{err} is an \kw{error}
  \item Pattern 2 limits the scope of \kw{x} and \kw{ok} to the if/else blocks
  \item More verbose than try/except, but every failure is handled where it happens
\end{itemize}

% ======================================================
\section{Loops: only \texttt{for}}
\begin{lstlisting}
for i := 0; i < 5; i++ {}   // C style
for n < 100 {}              // "while"
for {}                      // forever
for i, v := range items {}  // index, value
for _, v := range items {}  // like: for v in items
for k, v := range m {}      // key, value
for k := range m {}         // key only
for i := range 10 {}        // 0 to 9 (Go 1.22+)
\end{lstlisting}
\begin{itemize}
  \item No \kw{while} keyword; \kw{range} instead of \kw{in}
  \item Trap: \kw{for v := range items} gives \kw{v} the \textbf{index}, not the value
\end{itemize}

% ======================================================
\section{switch}
\begin{lstlisting}
switch n {
case 1:
    fmt.Println("one")
    fallthrough             // run the next case's body anyway
case 2, 3:                  // several values
    fmt.Println("one, two or three")
default:
    fmt.Println("other")
}                           // n == 1 prints both lines
\end{lstlisting}
\begin{itemize}
  \item No \kw{break} needed: the switch ends after the matching case
  \item \kw{fallthrough} skips the next case's condition check
  \item A \kw{case} can hold any code: an \kw{if}, a loop, even another \kw{switch}
  \item Cleaner than long if/else chains; \kw{gofmt} lines \kw{case} up with \kw{switch}
\end{itemize}

% ======================================================
\section{Pointers}
\begin{lstlisting}
x := 5
p := &x
*p = 10            // x is now 10
pt := &Point{X: 1}
pt.X = 5           // no -> needed

func newCounter() *int {
    n := 0
    return &n      // fine in Go
}
\end{lstlisting}
\begin{itemize}
  \item Same as C: \kw{\&x}, \kw{*p}; \kw{nil} is \Go's \kw{NULL}
  \item Different: no pointer arithmetic, garbage collected, returning \kw{\&local} is safe
\end{itemize}

% ======================================================
\section{defer}
\begin{lstlisting}
f, err := os.Open("config.txt")
if err != nil {
    return err
}
defer f.Close()    // runs when the function returns
\end{lstlisting}
\begin{itemize}
  \item Like Python's \kw{with}, but tied to the whole function, not a block
  \item Several defers run in reverse order
  \item C++ has no \kw{defer}; RAII destructors are the closest match
\end{itemize}

% ======================================================
\section{Small things that won't compile}
{\footnotesize
\begin{tabularx}{\linewidth}{@{}>{\ttfamily}p{0.40\linewidth}>{\ttfamily}X@{}}
\rowcolor{tablehead}\textsf{\textbf{Error}} & \textsf{\textbf{OK}}\\
if len(s) \{           & if len(s) > 0 \{\\
n + 1.5  \textcolor{comment}{// n is int} & float64(n) + 1.5\\
\} \textcolor{comment}{(newline)} else \{ & \} else \{\\
x = i++, ++i           & i++ \textsf{then} x = i\\
\end{tabularx}}
\begin{itemize}
  \item No truthiness, no implicit conversions, no prefix \kw{++}/\kw{--}
\end{itemize}

% ======================================================
\section{Functions}
\begin{lstlisting}
func add(a, b int) int {          // shared types grouped
    return a + b
}

func divide(a int, b int) (int, error) {
    if b == 0 {
        return 0, errors.New("cannot divide by zero")
    }
    return a / b, nil
}
result, err := divide(10, 2)
q, _ := divide(9, 3)              // _ discards a value
\end{lstlisting}
\begin{itemize}
  \item No prototypes or headers: a function can be anywhere in the package
  \item \kw{error} is always the last return value
  \item Failure: zero value + error. Success: result + \kw{nil}
  \item Multiple returns look like tuples, but \Go{} has no tuple type
\end{itemize}

% ======================================================
\section{Structs}
\begin{lstlisting}
type Person struct {
    Name string
    Age  int
}
p := Person{Name: "Ana", Age: 30}
r := Person{Name: "Ana"}   // Age 0
var q Person               // Name "", Age 0

inv := Inventory{}
inv.Items = make(map[string]int)  // before adding!
\end{lstlisting}
\begin{itemize}
  \item \kw{p.Name} works on a struct or a pointer to one
  \item Other packages (incl.\ \kw{encoding/json}) only see Capitalised fields; use a tag for lowercase JSON keys:
  \item Map fields need \kw{make} before adding, or you get a runtime panic
\end{itemize}
\begin{lstlisting}
type User struct {
    Name string `json:"name"`   // {"name":"Ann"}
    age  int                    // skipped by json.Marshal
}
\end{lstlisting}

% ======================================================
\section{Methods \& receivers}
\begin{lstlisting}
func (p Person) Greet() string {  // value: works on a copy
    return "hi, I'm " + p.Name
}
func (p *Person) Birthday() {     // pointer: changes original
    p.Age++
}
p.Birthday()                      // Go takes &p for you
\end{lstlisting}
\begin{itemize}
  \item Value receiver: method only reads (like a C++ \kw{const} member function)
  \item Pointer receiver: method changes fields, or struct is large
  \item Don't mix them on one type: if one needs a pointer, use pointers for all
  \item A method is a function with a receiver before its name; same idea as passing a struct vs a pointer in C
  \item Methods are what interfaces are built on (\Go's version of duck typing)
\end{itemize}

% ======================================================
\section{NewX functions}
\begin{lstlisting}
func NewPerson(name string) *Person {
    return &Person{Name: name}
}
p := NewPerson("Ana")
\end{lstlisting}
\begin{itemize}
  \item No constructors / \kw{\_\_init\_\_}: by convention a plain \kw{New}+type function
  \item Only needed when a struct needs setup; otherwise use a literal
\end{itemize}

\end{multicols*}
\end{document}
